\section{The sharp global error and its parity profile}\label{sec:globalsharpproof}

We prove Theorem~\ref{thm:global}. Retain $K=k-\epsilon=2m$, $x=\sqrt{K/N}$, $\lambda=(2/3)K^{3/2}/\sqrt N$, and the definitions in \eqref{eq:actionq}, \eqref{eq:globalformula}, and \eqref{eq:globalprofiledef}. The exact normalization \eqref{eq:qnormalize} cancels the permutation denominator:
\begin{equation}\label{eq:globalcancel}
1+\mathcal E_{N,k}
=\frac{Q_k/(\mathcal L_kH_k)}{e^{Nq(x)}(1+x)^{\epsilon-1/2}}.
\end{equation}
This identity holds throughout $2\le k\le N$.

\subsection{The compact critical expansion}

Fix $0<d<C<\infty$ and suppose $dN^{1/3}\le k\le CN^{1/3}$. Every estimate in this subsection is uniform in that range and in parity. In particular, $K\to\infty$ and $\lambda$ stays in a compact subset of $(0,\infty)$.

Let $J$ be Poisson of mean $\lambda$, conditioned to have parity $\epsilon$. Independently, let $V$ be Poisson of mean
\[
\eta=\frac{3K^2}{4N}=\frac{27\lambda^2}{16K}.
\]
Set $T=J-\epsilon$ and $S=3T/2+2V$, both nonnegative integers. For integer $s\ge0$, define
\[
w_\epsilon(s)=\frac{(m-\epsilon)_s}{m^s}
\quad(0\le s\le m-\epsilon),
\]
and assign value zero outside this range. The pair availability identities \eqref{eq:availabilityeven} and \eqref{eq:availabilityodd} show that the contribution from blocks of sizes two, three and four to $Q_k/(\mathcal L_kH_k)$ is exactly
\begin{equation}\label{eq:globalconditioned}
e^\eta\E w_\epsilon(S).
\end{equation}
Differentiating the even and odd exponential series gives
\begin{equation}\label{eq:globaltriplemoments}
\E T=p_\epsilon(\lambda),\qquad
\E T^2=\lambda^2+(1-2\epsilon)p_\epsilon(\lambda).
\end{equation}
All fixed moments of $T$ are uniformly bounded here. The same holds for $S$, since $\eta=O(K^{-1})$.

Put $A(s)=s(s-1+2\epsilon)/(2m)$. The elementary availability estimate is
\[
0\le e^{-A(s)}-w_\epsilon(s)\le4s^3/m^2.
\]
For $s\le m/2$, expand the logarithm of the product and use
$0\le-\log(1-u)-u\le u^2/[2(1-u)]$. For larger $s$, the difference is at most one and $4s^3/m^2>1$, including when the product is assigned zero. Thus unavailable pair counts are included in the bound. Since $|e^{-a}-1+a|\le a^2/2$ for $a\ge0$,
\[
\E w_\epsilon(S)=1-\E A(S)+O(K^{-2}).
\]
Terms containing $V$ contribute $O(K^{-2})$ to $\E A(S)$. Substituting \eqref{eq:globaltriplemoments} into the remaining terms gives
\begin{equation}\label{eq:globalavailabilitymean}
\E A(S)=\frac1K\left(\frac94\lambda^2
+\frac34(1-2\epsilon)p_\epsilon(\lambda)\right)+O(K^{-2}).
\end{equation}

All blocks of size at least five together have weight at most
\[
\sum_{j\ge5}N\frac{j-1}{j}x^j
\le\frac{Nx^5}{1-x}=O(K^{-2}).
\]
Removing their availability factors bounds their additional positive contribution. Projecting onto the required parity and dividing by $\cosh\lambda$ or $\sinh\lambda$ preserves the bound on this compact parameter interval. Their total contribution to $Q_k/(\mathcal L_kH_k)$ is therefore $O(K^{-2})$. With $e^\eta=1+\eta+O(K^{-2})$, equations \eqref{eq:globalconditioned} and \eqref{eq:globalavailabilitymean} yield
\begin{equation}\label{eq:globalQcompact}
\frac{Q_k}{\mathcal L_kH_k}
=1+\frac1K\left(-\frac9{16}\lambda^2
+\frac34(2\epsilon-1)p_\epsilon(\lambda)\right)+O(K^{-2}).
\end{equation}

Now $x=3\lambda/(2K)$, while $q(x)=-x^4/4+O(x^5)$ and $\log(1+x)=x+O(x^2)$. Hence
\begin{equation}\label{eq:globalAcompact}
e^{Nq(x)}(1+x)^{\epsilon-1/2}
=1+\frac1K\left(-\frac9{16}\lambda^2
+\frac34(2\epsilon-1)\lambda\right)+O(K^{-2}).
\end{equation}
Dividing \eqref{eq:globalQcompact} by \eqref{eq:globalAcompact} gives
\begin{equation}\label{eq:globalcompact}
\mathcal E_{N,k}
=\frac{3(2\epsilon-1)}{4K}
\bigl(p_\epsilon(\lambda)-\lambda\bigr)+O_{d,C}(K^{-2}).
\end{equation}
Because $K=k+O(1)$ and $\lambda-\ell(k/N^{1/3})=O_{d,C}(N^{-1/3})$, this proves \eqref{eq:globalprofile} uniformly on every fixed compact interval $d\le k/N^{1/3}\le C$.

\subsection{Two real saddles with uniform constants}

The exact identity
\begin{equation}\label{eq:globalrealintegral}
Q_k=\frac1{(N-k)!}\int_0^\infty e^{-v}v^{N-k}(N-v)^k\,dv
\end{equation}
follows by expanding $(N-v)^k$ and integrating each monomial. Splitting at $v=N$ and substituting $v=N(1\mp t)$ gives $Q_k=T_++(-1)^kT_-$, where
\begin{align*}
T_+&=\frac{N^{N+1}e^{-N}}{(N-k)!}
\int_0^1e^{Nt}t^k(1-t)^{N-k}\,dt,\\
T_-&=\frac{N^{N+1}e^{-N}}{(N-k)!}
\int_0^\infty e^{-Nt}t^k(1+t)^{N-k}\,dt.
\end{align*}
Both positive integrals have a unique maximum at $u=\sqrt{k/N}$. For integers $2\le k\le N/4$ and $k\to\infty$,
\begin{equation}\label{eq:globalrealLaplace}
T_\pm=B_\pm(1+O(k^{-1})),
\end{equation}
with an absolute implied constant and an absolute sufficient lower threshold for $k$, where
\begin{align}
B_+&=\frac{(Nu)^k}{\sqrt{2(1+u)}}
\exp\bigl(N[u-u^2-(1-u^2)\log(1+u)]\bigr),\label{eq:globalBplus}\\
B_-&=\frac{(Nu)^k}{\sqrt{2(1-u)}}
\exp\bigl(N[-u-u^2-(1-u^2)\log(1-u)]\bigr).\label{eq:globalBminus}
\end{align}

Here is the uniform Laplace justification. Rescale $t=us$, subtract the phase at $s=1$, and divide by $k$. The resulting phases satisfy
\[
F_\pm'(s)=\frac{1-s^2}{s(1\mp us)},\qquad
F_\pm''(1)=-\frac2{1\mp u}.
\]
Their derivatives through order six are uniformly bounded on a fixed neighborhood of one for $0\le u\le1/2$. Their second derivatives are bounded away from zero there. Outside this neighborhood, the left part has a fixed negative phase gap. On the right, $F_\pm'$ stays below a fixed negative constant. The right tails therefore have integrable exponential bounds uniform in their possibly unbounded endpoints.

In the central neighborhood, Taylor expansion through the quartic term and integration over a symmetric interval cancel the odd cubic contribution. The remaining absolute error is $O(k^{-3/2})$ against a leading Gaussian integral of order $k^{-1/2}$. The outer integrals are exponentially small. The second derivatives above give the Gaussian factors $\sqrt{\pi(1\mp u)/N}$ in the original variable. Applying Stirling's formula to $(N-k)!$ adds $O((N-k)^{-1})=O(k^{-1})$ and gives \eqref{eq:globalBplus} and \eqref{eq:globalBminus}. This proves \eqref{eq:globalrealLaplace}, including as $u\to0$.

Define the positive part of the target by
\begin{equation}\label{eq:globalP}
P_{N,k}=\frac{\mathcal L_k}{2\lambda^\epsilon}
e^{Nh(x)}(1+x)^{\epsilon-1/2}.
\end{equation}
Then
\begin{equation}\label{eq:globaltargetsplit}
\mathcal L_kH_ke^{Nq(x)}(1+x)^{\epsilon-1/2}
=P_{N,k}\bigl(1+(-1)^\epsilon e^{-2\lambda}\bigr).
\end{equation}
For integers $2\le k\le N/4$ and $K\to\infty$,
\begin{equation}\label{eq:globalpositivecompare}
B_+/P_{N,k}=1+O(K^{-1}),
\end{equation}
again with an absolute implied constant and an absolute sufficient lower threshold for $K$. For even $k=K$, use
\[
\mathcal L_K=\sqrt2(NK/e)^{K/2}(1+O(K^{-1}))
\]
in \eqref{eq:globalBplus} and \eqref{eq:globalP}. For odd $k=K+1$, the exact normalizers give
\[
P_{N,K+1}/P_{N,K}=(K+1)(1+x)/x.
\]
Regard the logarithm of \eqref{eq:globalBplus} as a smooth function $g(v)$ of its degree. Direct differentiation gives
\[
g'(v)=\log(Nu_v(1+u_v))-\frac1{4Nu_v(1+u_v)},
\qquad u_v=\sqrt{v/N}.
\]
Its variation on $K\le v\le K+1$ is $O(K^{-1})$. Integrating over that interval gives the same adjacent-degree ratio up to relative $O(K^{-1})$. This proves \eqref{eq:globalpositivecompare} in both parities.

\subsection{Matching the negative saddle and the upper tail}

The two leading saddle terms have ratio
\begin{equation}\label{eq:globalnegativeratio}
\frac{B_-}{B_+}
=\exp\bigl(\operatorname{artanh}u-N\Delta(u)\bigr),
\qquad
\Delta(u)=2u-(1-u^2)\log\frac{1+u}{1-u}.
\end{equation}
For $u\le1/2$, $\Delta(u)=4u^3/3+O(u^5)$ with an absolute remainder bound. Put $\lambda_k=(2/3)k^{3/2}/\sqrt N$. The exact odd shift satisfies $0\le\lambda_k-\lambda\le u$. Thus
\[
\frac{B_-}{B_+}=e^{-2\lambda}e^\theta,
\qquad |\theta|\le C(u+Nu^5).
\]
When $k\le\sqrt N$, the right side tends uniformly to zero, and hence
\[
\left|\frac{B_-}{B_+}-e^{-2\lambda}\right|
\le C(u+Nu^5)e^{-2\lambda}.
\]
Combine this bound with \eqref{eq:globalrealLaplace}, \eqref{eq:globaltargetsplit}, and \eqref{eq:globalpositivecompare}. If $\lambda\ge1$, the denominator $1+(-1)^\epsilon e^{-2\lambda}$ is bounded away from zero. We obtain
\begin{equation}\label{eq:globalmiddlebound}
|\mathcal E_{N,k}|\le C\left[
K^{-1}+\left(\sqrt{k/N}+\frac{k^{5/2}}{N^{3/2}}\right)e^{-2\lambda}\right].
\end{equation}
There are absolute constants $C,N_0,K_0$ for which this holds for every integer $N\ge N_0$ and $2\le k\le\sqrt N$ satisfying $K\ge K_0$ and $\lambda\ge1$.

Write $c_k=k/N^{1/3}$. Since $K\ge k/2$,
$\lambda\ge c_k^{3/2}/(3\sqrt2)$. For fixed $C_0\ge4$, \eqref{eq:globalmiddlebound} implies
\begin{align}
&\limsup_{N\to\infty}N^{1/3}
\sup_{C_0N^{1/3}\le k\le\sqrt N}|\mathcal E_{N,k}|\notag\\
&\qquad\le C_1\left[
\frac2{C_0}+\sup_{t\ge C_0}(\sqrt t+t^{5/2})
e^{-2t^{3/2}/(3\sqrt2)}\right].\label{eq:globalmiddleupper}
\end{align}
Indeed, $N^{1/3}\sqrt{k/N}=\sqrt{c_k}$ and
$N^{1/3}k^{5/2}/N^{3/2}=c_k^{5/2}N^{-1/3}\le c_k^{5/2}$. The right side of \eqref{eq:globalmiddleupper} tends to zero as $C_0\to\infty$.

For the remaining upper orders, \eqref{eq:normalizedsaddle} divided by the exact hyperbolic factor gives
\[
|\mathcal E_{N,k}|\le C(K^{-1}+e^{-\lambda/8})
\qquad(\sqrt N\le k\le N)
\]
for sufficiently large $N$. Its right side is uniformly $O(N^{-1/2})$. Therefore
\begin{equation}\label{eq:globaltailinfty}
\lim_{C_0\to\infty}\limsup_{N\to\infty}
N^{1/3}\sup_{C_0N^{1/3}\le k\le N}|\mathcal E_{N,k}|=0.
\end{equation}
This includes the endpoint $k=N$.

\subsection{Fixed orders and the lower tail}

The identities $q(0)=0$ and $q'(x)=2x(\log(1+x)-x)$ give
$-x^4/4\le q(x)\le0$ for $x\ge0$. Also
$0\le\log D_k\le k^2/N$ for $k\le N/4$. It follows that
\[
\left|\log(\mathcal A_{N,k}/H_k)\right|
\le\frac54\frac{k^2}N+\frac12\sqrt{k/N}.
\]
The finite logarithmic comparison \eqref{eq:logerror} therefore yields
\[
\left|\log(1+\mathcal E_{N,k})\right|
\le6k^2/N+\tfrac12\sqrt{k/N}
\]
when $k\le N/4$ and $4k^2\le N$. For each fixed $d>0$, these conditions hold for every $2\le k\le dN^{1/3}$ once $N$ is sufficiently large. Exponentiating gives
\begin{equation}\label{eq:globaltailzero}
\limsup_{N\to\infty}N^{1/3}
\sup_{2\le k\le dN^{1/3}}|\mathcal E_{N,k}|
\le6d^2+\tfrac12\sqrt d.
\end{equation}
This tends to zero as $d\downarrow0$ and includes all fixed orders.

\subsection{The full profile and its maximum}

The inequalities $0\le p_\epsilon(t)\le t$ give
$|\Phi_\epsilon(c)|\le\sqrt c/2$. More explicitly,
\begin{align*}
\Phi_0(c)&=\frac3{4c}\frac{2\ell(c)}{e^{2\ell(c)}+1}>0,\\
\Phi_1(c)&=-\frac3{4c}
\left(1-\frac{2\ell(c)}{e^{2\ell(c)}-1}\right)<0.
\end{align*}
Both are bounded in absolute value by $3/(4c)$ and tend to zero as $c\to\infty$.

Split the full range at $dN^{1/3}$ and $C_0N^{1/3}$. On the middle range, \eqref{eq:globalcompact} gives uniform convergence to the profile. Equations \eqref{eq:globaltailzero} and \eqref{eq:globaltailinfty}, together with the profile bounds just proved, control the lower and upper ranges. First take $N\to\infty$, then $d\downarrow0$ and $C_0\to\infty$. This proves \eqref{eq:globalprofile} uniformly over every $2\le k\le N$.

For $c>0$,
\[
|\Phi_1(c)|-\Phi_0(c)
=\frac3{4c}\left(1-\frac{2\ell(c)}{\sinh(2\ell(c))}\right)>0.
\]
Thus the odd absolute profile determines the maximum in \eqref{eq:globalsharpconstant}. It is positive and tends to zero at both ends of $(0,\infty)$, so its maximum is attained. The uniform profile gives the upper bound in \eqref{eq:globalerror}. Taking an odd integer nearest a maximizing $c_*N^{1/3}$ gives the matching lower bound and the negative sign of the error.

It remains to prove uniqueness of the maximizing scale. Set $z=2\ell(c)=(4/3)c^{3/2}$. Apart from the positive factor $(3/4)^{1/3}$, the function to maximize is
\[
g(z)=z^{-2/3}\left(1-\frac z{e^z-1}\right).
\]
The derivative has the sign of $-\mathcal H(z)$, where
\[
\mathcal H(z)=2e^{2z}-(3z^2-z+4)e^z-z+2.
\]
Direct expansion gives
\[
\mathcal H(z)=-\frac{z^3}2-\frac{z^4}6
+\sum_{n\ge5}\frac{2^{n+1}-3n^2+4n-4}{n!}z^n.
\]
Each numerator in the sum is positive. It equals five at $n=5$, and its value at $n+1$ minus twice its value at $n$ is $3n^2-10n+5>0$ for $n\ge5$. Therefore $\mathcal H(z)/z^4$ is strictly increasing for $z>0$, from negative infinity to positive infinity. It has exactly one positive zero. The function $g$ increases before this zero and decreases afterward, proving \eqref{eq:globalmaxequation} and uniqueness. This completes the proof of Theorem~\ref{thm:global}.
